\begin{helper}
Let $\nu$ be any measure on activation--priority pairs $(A,\pi)$, and use the
output map from \noderef{RandomIndependentSetSampling}: for such a pair write
\[
I(A,\pi)=\{z\in A:\pi(w)<\pi(z)\text{ for every }w\in A\cap N_G(z)\}.
\]
For every fixed finite vertex set $Q$,
\[
\nu\{(A,\pi):Q\subseteq I(A,\pi)\}
=\sum_{A:Q\subseteq A}
\nu\{(A',\pi):A'=A\text{ and every }q\in Q\text{ beats every activated
neighbor in }A\}.
\]
Moreover, for each activation set $A$ and priority function $\pi$,
\[
Q\subseteq I(A,\pi)
\quad\Longleftrightarrow\quad
Q\subseteq A\text{ and every }q\in Q\text{ beats every activated neighbor
in }A.
\]
\end{helper}
\begin{proof}
The activation coordinate is discrete and the priority coordinates have
their Borel product sigma algebra. Thus the activation equality $A'=A$
is measurable, as is every strict comparison $\pi(w)<\pi(q)$.
Each fiber displayed in the sum is a finite intersection of such events,
and therefore is measurable.

The pointwise characterization is obtained by applying
\noderef{SamplingOneVertexOutputMembership} to each vertex $q\in Q$.
Indeed, $q\in I(A,\pi)$ is equivalent to the conjunction that $q\in A$ and
that every activated neighbor $w\in A$ of $q$ has $\pi(w)<\pi(q)$.  Thus
$Q\subseteq I(A,\pi)$ holds exactly when this conjunction holds for every
$q\in Q$, which is the displayed characterization.

Now partition the event $\{(A,\pi):Q\subseteq I(A,\pi)\}$ according to the
first coordinate $A$.  By the characterization just proved, only activation
sets satisfying $Q\subseteq A$ can occur, and for such an $A$ the fiber is
precisely the event that the first coordinate equals $A$ and every
$q\in Q$ beats every activated neighbor in $A$.  These fibers are pairwise
disjoint, because a single sample has only one activation set.  Since the
vertex set is finite, there are only finitely many activation sets, so finite
additivity of $\nu$ over this disjoint union gives the stated sum.
\end{proof}
